\documentclass[english, Lau, oneside]{sapthesis}
\usepackage[utf8]{inputenc}
\usepackage{indentfirst}
\usepackage{microtype}
\usepackage{lettrine}
\linespread{0.9}
\usepackage[nottoc, notlof, notlot]{tocbibind}
\usepackage{hyperref}
\hypersetup{
hyperfootnotes=true,
bookmarks=true,
colorlinks=true,
linkcolor=red,
linktoc=page,
anchorcolor=black,
citecolor=red,
urlcolor=blue,
pdftitle={AI-Based Visual Monitoring System for Real-Time Anomaly Detection in Large-Scale 3D Printing},
pdfauthor={Aleksandar Kirilov},
pdfkeywords={3D printing, anomaly detection, AI, FDM}
}
\title{AI-Based Visual Monitoring System for Real-Time Anomaly Detection in Large-Scale 3D Printing}
\author{Aleksandar Kirilov}
\IDnumber{2059920}
\course{Applied Computer Science and Artificial Intelligence}
\courseorganizer{Facoltà di Ingegneria dell'informazione, informatica e statistica}
\submitdate{2025/2026}
\copyyear{2025}
\advisor{Prof. Marco Raoul Marini}
% TODO: add real co-advisor name, or delete this line if there is none.
% \coadvisor{Dr. ...}
\authoremail{kirilov.2059920@studenti.uniroma1.it}
% TODO: set your real graduation/exam date before submission (e.g. October 2026).
\examdate{TODO: exam date}
% TODO: replace with the real examination commission members before submission.
\examiner{Prof. ...} \examiner{Prof. ...} \examiner{Prof. ...} \examiner{Prof. ...} \examiner{Prof. ...} \examiner{Prof. ...} \examiner{Prof. ...}
\begin{document}
\frontmatter
\maketitle

\begin{abstract}
Large-scale Fused Deposition Modeling (FDM) printers enable the production of large components such as humanoid robot parts but are highly susceptible to print failures that can result in significant material and time waste. This thesis presents the design and development of an AI-based visual monitoring system for the real-time detection of anomalies during the printing process. A compact camera is installed inside the printer to continuously monitor the printed surface. The system processes the video stream using computer vision and machine learning techniques to identify printing defects such as stringing, warping, layer shift, under-extrusion, and the "blob of death". The main objective of this work is to create a low-cost, scalable, and autonomous solution that can detect anomalies early and notify the user or pause the printing process, improving reliability and reducing waste in large-scale 3D printing environments.
\end{abstract}

\tableofcontents

\mainmatter

% --------------------------------------------------
\chapter{Introduction}
\label{chap:introduction}

\lettrine[lines=3]{L}{arge}-scale 3D printing is increasingly used in academic and industrial environments to manufacture large mechanical parts, including structural elements for humanoid robots. Despite its potential, the process remains vulnerable to a wide range of printing anomalies that often lead to failed prints and significant material loss.

The goal of this thesis is to design and implement an AI-based visual monitoring system capable of detecting common anomalies in real time during the printing process of a large-scale FDM printer. Unlike traditional monitoring systems that rely on sensors or post-process inspection, the proposed solution focuses on visual inspection using a low-cost camera combined with computer vision and machine learning techniques.

\section{Motivation}
Large prints can take several hours or even days to complete. An undetected failure may result in wasted material, machine time, and energy consumption. Early detection of such failures is therefore essential in research environments where resources are limited.

\section{Objectives}
The main objectives of this thesis are:
\begin{itemize}
\item Identify and study the most common anomalies in large-scale FDM printing.
\item Design a low-cost visual monitoring system using a compact camera and an edge computing device.
\item Develop an AI-based model capable of detecting anomalies in real time.
\item Integrate the system with a 3D printer and test its performance.
\item Evaluate the effectiveness of the system in terms of accuracy and response time.
\end{itemize}

\section{Thesis Structure}
This thesis is organized as follows:
\begin{itemize}
\item Chapter 2 presents the fundamentals of FDM printing and the targeted defect types.
\item Chapter 3 reviews existing approaches for defect detection and monitoring.
\item Chapter 4 describes the overall system architecture and design.
\item Chapter 5 details the dataset creation, preprocessing, and synthetic data generation.
\item Chapter 6 covers the model architecture, training protocol, and validation results.
\item Chapter 7 presents the real-time system implementation.
\item Chapter 8 discusses the experimental evaluation and real-world testing.
\item Chapter 9 concludes the work and proposes future developments.
\end{itemize}

% --------------------------------------------------
\chapter{FDM Printing and Defect Types}
\label{chap:defects}

\section{FDM Printing Process}
Fused Deposition Modeling (FDM) is one of the most widely adopted Additive Manufacturing (AM) technologies, especially in research and prototyping contexts. The working principle is straightforward: a thermoplastic filament --- commonly PLA, ABS, or PETG --- is fed into a heated nozzle, where it melts and is deposited onto a build surface in a precisely controlled path. This process is repeated layer by layer, guided by a G-code file generated from a 3D model by slicing software, until the full three-dimensional geometry is constructed.

The main hardware components involved in the process are:
\begin{itemize}
    \item \textbf{Extruder} --- feeds the filament at a controlled rate into the hot end.
    \item \textbf{Hot End / Nozzle} --- heats the filament to between 180°C and 260°C depending on the material.
    \item \textbf{Heated Build Plate} --- maintains a surface temperature to improve first-layer adhesion.
    \item \textbf{Motion System} --- stepper motors control the precise X, Y, and Z movement of the toolhead.
\end{itemize}

Print quality depends heavily on a set of interdependent parameters: nozzle temperature, bed temperature, print speed, layer height, cooling fan speed, and retraction settings. Small deviations in any of these --- from incorrect calibration, environmental factors, or hardware wear --- can produce visible defects that compromise the geometric accuracy and structural integrity of the printed part.

\section{Challenges in Large-Scale Printing}
Scaling FDM to large-format printers introduces challenges that do not exist, or exist to a much lesser degree, at desktop scale.

\textbf{Thermal management:} As the printed surface grows, maintaining uniform temperature becomes harder. Uneven cooling creates internal stresses that cause the material to shrink and lift off the build plate. This is one of the most common failure modes in large-format printing --- even a small temperature gradient across a 500\,mm or larger bed can generate enough stress to delaminate the first layers.

\textbf{Temporal exposure:} Large prints can require 20 to 48 hours of continuous operation. A defect that appears at hour 10 of a 30-hour print can go entirely unnoticed until the operator checks the machine, at which point the damage is irreversible. This makes real-time automated monitoring particularly valuable \cite{zhang2025}.

\textbf{Mechanical precision over time:} Long print runs put stress on the motion system. Stepper motors that skip steps, loose belts, or external vibrations can cause a sudden misalignment of the toolhead mid-print, resulting in misaligned layers across the full height of the part.

\textbf{Material and resource cost:} Large prints consume significantly more filament, electricity, and machine time. A failed large-format print represents a meaningful cost in materials and availability, particularly in academic lab environments where resources are limited.

\section{Common Anomalies}
A wide range of defect types can occur during FDM printing \cite{armin2025}. This section describes the defect categories targeted in this work, selected based on their visual detectability from a camera mounted inside the printer enclosure. Five of them --- Spaghetti, Warping, Layer Shifting, Stringing, and Cracking --- are addressed by the detection model presented in this thesis; the remaining two, Off-Platform and Blob of Death, belong to the system's intended defect taxonomy and are targeted for a future training iteration once sufficient data has been collected.

\textbf{Spaghetti} occurs when the nozzle loses adhesion to the print and continues extruding filament into the air, producing a chaotic tangle of strands. Once it begins, the print is unrecoverable and the deposited material can jam the motion system or foul the nozzle.

\textbf{Warping} is the upward curling of the bottom edges of a print caused by differential thermal contraction during cooling. On large prints, the lifting force can be strong enough to fully detach the part from the bed. It is most visible in the early layers and tends to worsen progressively if not caught early.

\textbf{Layer Shifting} occurs when the print head loses its reference position mid-print, causing all subsequent layers to be offset in the X or Y direction. This produces a characteristic staircase effect on the exterior of the part. Common causes include belt slippage, mechanical obstruction, or stepper motor overload.

\textbf{Stringing} refers to thin filament threads left between separate features of a print as the nozzle travels across open gaps. The nozzle leaks a small amount of molten material during non-print moves, leaving hair-like strings behind. While not always catastrophic, heavy stringing indicates incorrect retraction settings and degrades the surface quality of the final part.

\textbf{Off-Platform} describes the condition where the printed part has fully detached from the build plate and the nozzle is depositing material onto nothing. It is a terminal failure state, often coinciding with spaghetti generation, and is visually identifiable by the absence of a coherent structure on the bed.

\textbf{Cracking} occurs when layers delaminate from each other, typically due to insufficient interlayer bonding. This can result from printing too fast, at too low a temperature, or in a cold environment that causes rapid solidification before adequate adhesion forms between layers.

\textbf{Blob of Death} is a large accumulation of molten filament that builds up around the nozzle, typically after debris or a detached first layer causes extruded material to adhere to the hot end rather than to the print. The blob grows layer after layer and can fully encase the nozzle, jamming the toolhead and, in severe cases, damaging the hot end. It is included as a target defect because it is both visually distinctive and mechanically dangerous, but it is reserved for a future training iteration rather than the current model.

Other defect types exist in the literature --- such as under-extrusion, over-extrusion, and Z-banding --- but their visual signatures are either more subtle or require higher camera resolution to detect reliably, and are therefore outside the scope of this work.

% --------------------------------------------------
\chapter{Related Work and State of the Art}
\label{chap:related}

The automated detection of defects in additive manufacturing has been approached from several complementary directions, ranging from hardware sensing to purely software-based visual inspection. This chapter surveys the main families of approaches, broadly following the taxonomy proposed in recent reviews of the field \cite{armin2025}, and identifies the gaps that motivate the design choices of this work.

\section{Traditional Monitoring Approaches}
The earliest and still most widespread monitoring strategies rely either on physical sensors or on human supervision.

\textbf{Sensor-based monitoring} instruments the printer with dedicated hardware: thermocouples and thermistors track nozzle and bed temperature, accelerometers and current sensors detect abnormal motion or stepper-motor load, and acoustic-emission sensors capture the characteristic sound signatures of extrusion faults. Zhang et al.\ \cite{zhang2025} demonstrate that fusing such process signals with a data-driven model enables in-situ anomaly detection and even closed-loop parameter optimisation. While effective for the failure modes they are designed to capture, sensor-based methods share two limitations: they require additional hardware to be integrated into the machine, increasing cost and complexity, and most process signals are global, providing little information about \textit{where} on the bed a defect has occurred.

\textbf{Manual visual inspection} remains the default in many laboratory and prototyping settings: an operator periodically checks the print, either in person or through a webcam stream. This requires no specialised hardware but does not scale to the long, unattended print runs typical of large-format FDM, where a failure at hour ten of a thirty-hour print can go unnoticed until it is irreversible.

\textbf{Post-process inspection} techniques --- such as coordinate-measuring machines, X-ray computed tomography, and dimensional metrology --- offer high-accuracy characterisation of the finished part but are inherently retrospective: they confirm that a print has failed only after the material and machine time have already been spent. None of these traditional approaches satisfies the requirement of early, autonomous, in-situ detection on large prints.

\section{Computer Vision Techniques}
Camera-based monitoring removes the need for additional process sensors and provides direct spatial information about the print. Classical computer-vision pipelines approach defect detection through hand-engineered image processing \cite{armin2025}. Common strategies include background subtraction and frame differencing to isolate the printed region from the bed, edge- and contour-detection to measure geometric deviations, optical-flow analysis to detect unexpected motion such as detached filament, and texture descriptors to characterise surface quality. A particularly influential idea is \textit{layer-wise reference comparison}, in which each printed layer is compared against an expected image rendered from the sliced G-code, and significant discrepancies are flagged as anomalies.

These techniques are computationally light and interpretable, which makes them attractive for embedded deployment. However, they are notoriously brittle in practice: their performance depends heavily on controlled and consistent lighting, fixed camera geometry, and careful per-setup calibration. A change in filament colour, a shift in ambient illumination, or a slightly different camera angle can invalidate hand-tuned thresholds, which limits their robustness in the variable conditions of a real printer enclosure.

\section{Deep Learning Approaches}
The limitations of hand-engineered features have driven a decisive shift toward deep learning, in which convolutional neural networks (CNNs) learn discriminative features directly from data. The effectiveness of deep CNNs has been established across a broad range of visual recognition tasks, from human action classification \cite{cascio2017} to body-affect recognition \cite{fagioli2019}, demonstrating that learned hierarchical features generalise across visual domains far better than hand-crafted descriptors. This same property underlies their success in defect detection.

Within additive manufacturing, deep-learning methods fall into three groups. \textbf{Image classification} models answer the binary or multi-class question of \textit{whether} a given frame contains a defect, but provide no spatial localisation. \textbf{Object detection} models, such as the YOLO (You Only Look Once) family, instead predict bounding boxes, identifying \textit{what} defect is present and \textit{where} --- information that is essential for a system intended to localise and react to failures. \textbf{Semantic segmentation} offers the finest spatial granularity, at a higher computational cost. Lightweight backbones such as EfficientNet, and more recently transformer-based detectors, have also been explored to balance accuracy against inference cost.

The work most closely related to this thesis is that of Cenedese \cite{cenedese2022}, who implemented a 3D-printing anomaly-detection system combining YOLOv5 for localisation with an EfficientNet-Lite classifier. The present work follows the same detection-based philosophy but targets a more recent architecture (YOLOv8s), a different deployment scenario --- a global, frame-mounted view of a large-format bed rather than desktop close-ups --- and closed-loop actuation through the printer's control API.

\section{Research Gaps}
Three gaps emerge from the literature and directly motivate the approach taken in this thesis.

First, \textbf{a domain gap in the available data}. As discussed in Section~\ref{sec:synthetic_data}, the publicly available defect datasets consist overwhelmingly of macro close-ups captured for cataloguing purposes, whereas a frame-mounted monitoring camera observes the entire bed from a distance. Models trained directly on close-ups therefore generalise poorly to the global field of view, which motivates the synthetic composition pipeline developed in this work.

Second, \textbf{a focus on desktop rather than large-format printers}. Most prior systems are demonstrated on small desktop machines, where the consequences of a failed print are minor. The specific challenges of large-format printing --- long unattended runs, severe thermal gradients, and high material cost --- have received comparatively little attention, despite being precisely where automated monitoring delivers the greatest value.

Third, \textbf{a gap between detection and action}. Many works stop at reporting detection accuracy and do not close the loop by intervening in the print. A practical monitoring system must not only detect a defect reliably but also run on constrained edge hardware and trigger a corrective action --- pausing the print --- while tolerating the false positives that are inevitable in real conditions. Addressing this end-to-end requirement, from a low-cost camera to an automatic pause command, is the central objective of this work.

% --------------------------------------------------
\chapter{System Design and Architecture}
\label{chap:architecture}

This chapter describes the overall design of the monitoring system: the physical hardware on which it runs, the modular software architecture that connects its components, and the operational workflow that ties them together during a print. The implementation details of each software component are deferred to Chapter~\ref{chap:implementation}.

\section{Hardware Design}
The system is designed as a self-contained, low-cost appliance that can be attached to a printer without modifying the machine itself.

\textbf{Edge device.} All computation runs on a \textbf{Raspberry Pi} single-board computer. This choice is central to the design philosophy: the Raspberry Pi is inexpensive, low-power, and self-contained, allowing inference to run \textit{on the device}, next to the printer, with no dependence on a cloud service or a powerful host workstation. Running at the edge keeps the system private, removes network latency from the detection loop, and makes it inexpensive to replicate across a fleet of machines.

\textbf{Camera.} A compact camera is mounted inside the printer enclosure with a global view of the print bed, so that the entire printed surface is captured in a single frame. The software is deliberately camera-agnostic: it accepts either the Raspberry Pi camera module connected over the CSI interface or a standard USB webcam, both exposed through the same acquisition interface (Chapter~\ref{chap:implementation}). Stable, consistent illumination inside the enclosure is important, since the model is sensitive to large changes in lighting; the glow of the heated bed and any enclosure lighting are therefore part of the operating environment the camera must tolerate.

\textbf{Target printer.} The deployment target is an \textbf{Elegoo OrangeStorm Giga}, a large-format FDM printer with a build volume on the order of $800 \times 800 \times 1000$\,mm, upgraded in this project to produce structural parts for a humanoid robot. This is precisely the high-stakes scenario that motivates the work (Chapter~\ref{chap:defects}): such parts require prints lasting many hours or days and consume large quantities of filament, so an undetected failure is extremely costly in both material and machine time. The very large bed also means that defects can appear far from the centre of the frame and that an individual defect occupies only a small fraction of the field of view --- two properties that directly shaped the data and model design described in Chapters~\ref{chap:model} and~\ref{chap:model_development}.

\textbf{Connectivity.} The Raspberry Pi communicates with the printer's controller --- an OctoPrint or Moonraker server running on or alongside the machine --- over the local network. This link carries the status queries and pause commands that allow the monitoring system to act on the print without any physical modification to the printer's electronics.

\section{Software Architecture}
The software is organised as a linear pipeline of decoupled stages, each with a single responsibility. This modularity is what allows the same control logic to run unchanged against either recorded video or live hardware (Chapter~\ref{chap:implementation}).

\begin{enumerate}
    \item \textbf{Data acquisition} --- a camera abstraction grabs frames from the live camera (or, for testing, a recorded video) at a fixed sampling rate.
    \item \textbf{Preprocessing} --- each frame is resized and letterboxed to the $640 \times 640$ input expected by the network (Section~\ref{sec:dataset}).
    \item \textbf{Anomaly detection} --- the YOLOv8s model performs detection and localisation. In a deep convolutional detector, feature extraction is not a separate hand-engineered stage but is learned internally by the network, which is one of the principal advantages of the deep-learning approach over the classical pipelines surveyed in Chapter~\ref{chap:related}.
    \item \textbf{Decision filtering} --- per-class confidence thresholds and a temporal persistence filter convert noisy per-frame detections into a stable decision, suppressing isolated false positives.
    \item \textbf{Response and notification} --- when a defect is confirmed, the printer-control component pauses the print and a notification component alerts the operator.
\end{enumerate}

Each stage corresponds to a distinct, independently testable software component, and the interfaces between them are narrow enough that any single stage --- the camera source, the model weights, or the printer endpoint --- can be replaced without disturbing the rest of the system.

\section{Operational Workflow}
During a print, the system runs the following loop continuously:

\begin{enumerate}
    \item A frame is captured from the camera once per second.
    \item The frame is preprocessed and passed to the detector, which returns any defects it finds together with their class and confidence.
    \item Detections below the per-class confidence threshold are discarded, and the remainder update the temporal persistence state.
    \item If no defect persists, the loop simply continues to the next frame.
    \item If a defect persists across the required number of consecutive frames, the system treats it as confirmed: it verifies that a print is genuinely in progress, issues a pause command to the printer, alerts the operator, and resets its persistence state so as not to re-trigger on the now-frozen defect.
    \item The operator can then inspect the print and decide whether to resume or cancel the job.
\end{enumerate}

This workflow runs unattended for the full duration of a print, turning a multi-day, high-value job on the large-format machine from an unsupervised risk into a continuously monitored process.

% --------------------------------------------------
\chapter{Dataset and AI Model Development}
\label{chap:model}

\section{Dataset Creation}
The single greatest obstacle to training a visual monitoring system for large-format FDM is the scarcity of suitable data. As established in Chapter~\ref{chap:related}, the publicly available defect datasets are dominated by macro close-ups that lack the spatial context of a frame-mounted camera, and no public corpus captures the specific global view of the target large-format bed. The dataset for this work was therefore assembled as a \textit{hybrid corpus} drawn from several complementary sources, each chosen to supply a property that the others lack.

\textbf{Real defect imagery.} The primary source of genuine defect appearance is the publicly available ``3D Printing Defect Dataset'' on Kaggle, which provides close-up images for the five trained classes. These images serve two distinct roles: they are the foreground material for the synthetic composition pipeline (Section~\ref{sec:synthetic_data}), and a subset is also injected directly into training as real samples, so that the model is exposed to authentic defect texture rather than to composites alone.

\textbf{Supplementary defect imagery.} The Kaggle close-ups do not cover every viewing condition the deployed system will encounter. In particular, stringing observed from a shallow side angle was poorly represented, so a small set of wide-angle stringing images was gathered from public hobbyist sources to broaden that class.

\textbf{Clean printer-bed backgrounds.} A set of photographs of the target large-format printer bed was captured to act as backgrounds for synthetic composition. Grounding the synthetic samples in the real deployment environment teaches the model the specific lighting, texture, and geometry of the bed it will ultimately monitor.

\textbf{Negative (defect-free) samples.} Teaching the model what a healthy print looks like is as important as teaching it the defects. Three negative sources were combined: photographs of successful prints on the target machine, a sample of clean webcam frames drawn from a public ``no-defect'' set, and a small, growing set of \textit{hard negatives} --- defect-free images that the model had previously misclassified, fed back into training to suppress specific recurring false positives.

Each real defect close-up is annotated with a single bounding box covering the central region of the frame, since the defect fills most of a close-up shot; synthetic composites are annotated automatically from the known paste location and size; and negative images carry empty annotations, signalling that no object is present. The way these sources are balanced and combined into the final training corpus --- including the deliberate decision to keep the colour-biased Warping and Cracking classes synthetic-only --- is detailed in Section~\ref{sec:dataset}, and the rationale for that decision is analysed alongside the real-world results in Section~\ref{sec:results}.

\section{Data Preprocessing}
Before training, the heterogeneous raw images from these sources are normalised into a single consistent format.

\textbf{Resolution.} All samples are standardised to the $640 \times 640$ pixel input expected by YOLOv8. Background photographs are cropped to a square aspect ratio and resized so that no part of the bed is geometrically distorted, and synthetic composites are generated directly at this resolution (Section~\ref{sec:synthetic_data}). At inference time, incoming frames are letterboxed to the same dimensions, preserving aspect ratio so that the geometry seen during training matches the geometry seen in deployment.

\textbf{Label encoding.} Annotations follow the standard YOLO format: one text file per image, each line listing a class index followed by the bounding-box centre and size, normalised to the $[0,1]$ range. Synthetic composites receive exact boxes derived from the compositing step; real close-ups receive the fixed near-full-frame box described above; and negative images are paired with empty label files, which YOLOv8 interprets as ``background only'' and uses to penalise spurious detections.

\textbf{Normalisation.} Pixel intensities are scaled to the $[0,1]$ range as part of the standard YOLOv8 input pipeline, applied uniformly to every frame so that the statistics seen during training and inference are identical.

\textbf{Train/validation split.} The complete corpus is partitioned once into an 80/20 train/validation split. The split spans every class as well as the negative samples, ensuring both that each class is represented in the validation set and that no image leaks between the two partitions. Keeping the split fixed and held out is what makes the validation metrics reported in Section~\ref{sec:results} a trustworthy estimate of generalisation.

\textbf{Runtime augmentation.} The preprocessing described here produces the static, on-disk dataset. A second, stochastic layer of transformation --- mosaic and mix-up composition, geometric perturbation, and HSV colour jitter --- is applied on the fly during training to expand the effective diversity of the data and to simulate the optical variability of a real camera. Because these augmentations are a property of the training protocol rather than of the stored dataset, they are described in detail in Section~\ref{sec:training}.

\section{Synthetic Data Generation and Domain Adaptation}
\label{sec:synthetic_data}

A significant challenge in developing visual monitoring systems for large-scale Additive Manufacturing (AM) is the scarcity of publicly available datasets that match the specific field of view (FoV) of a frame-mounted camera. While public repositories provide ample examples of defects such as ``Spaghetti'' (filament entanglement) or ``Layer Shifting,'' these are predominantly captured via macro-photography, lacking the spatial context of the full print bed. To address this \textit{domain gap}, a synthetic data generation pipeline was developed to project high-fidelity defect morphologies onto global-view backgrounds.

\subsection{Methodology}
The generation pipeline functions as a semantic composition engine, combining foreground defect features with varied background environments. The process consists of three distinct stages:

\begin{enumerate}
    \item \textbf{Morphological Extraction:} 
    Raw defect samples were sourced from the ``3D Printing Defect Dataset'' (Kaggle). To isolate the defect features from their original noisy backgrounds, a U-2 Net based background removal tool (`rembg`) was utilized. This generated precise alpha masks for defect classes including warping, stringing, and delamination, effectively turning them into independent visual objects.

    \item \textbf{Background Normalization:}
    A specific set of background images capturing the target large-scale printer bed was collected. These images were pre-processed and cropped to a square aspect ratio to align with the standard input tensor dimensions of the YOLOv8 architecture ($640 \times 640$ pixels). This step ensures that the model learns the specific lighting and texture of the intended deployment environment.

    \item \textbf{Centroid Alignment and Adaptive Scaling:}
    Unlike traditional domain randomization, which places objects stochastically, this research implemented a \textit{Centroid Alignment} strategy. The extracted defect features are dynamically resized to occupy between 50\% and 80\% of the print bed area and are composited at the geometric center of the background. 
\end{enumerate}

\subsection{Justification of the Approach}
This synthetic approach serves two critical functions for the reliability of the monitoring system:

\begin{itemize}
    \item \textbf{Feature Scale Invariance:} By resizing macro-defects to fit within the global view, the model is trained to recognize anomalies based on their morphological structure rather than their relative screen size.
    \item \textbf{Region of Interest (ROI) Focus:} By enforcing centroid alignment, the training process implicitly guides the neural network's attention mechanisms toward the print bed area, reducing the likelihood of False Positives triggered by objects in the surrounding environment (e.g., tools or cabling outside the printer frame).
\end{itemize}

The synthetic generator produces on the order of 150 labelled composites per class, providing the spatially-diverse, global-view component of the hybrid dataset described in Section~\ref{sec:dataset}. As detailed there, these synthetic samples are combined with real defect close-ups and clean-print negatives to form the final training corpus used for the deployed model.

\chapter{Model Development and Experimental Analysis}
\label{chap:model_development}

This chapter details the development of the computer vision-based anomaly detection system. It covers the architectural selection of the neural network, the engineering of the hybrid dataset, the training protocol employed, and a critical analysis of the experimental results obtained through both synthetic validation and real-world testing.

\section{Deep Learning Framework}
\label{sec:framework}

To address the constraints of edge deployment, specifically the limited computational resources of a Raspberry Pi, while maintaining high detection accuracy, the \textbf{YOLOv8s (Small)} architecture was selected. This model represents a meaningful evolution from previous iterations (YOLOv5 and YOLOv7) through key architectural refinements that improve both accuracy and gradient flow efficiency.

\subsection{Network Architecture}
YOLOv8 replaces the traditional C3 bottleneck modules with \textbf{C2f (Cross-Stage Partial Bottleneck with two convolutions)} modules. The C2f module enables richer gradient flow by combining high-level semantic features with low-level spatial features without significantly increasing computational cost.

\begin{equation}
    \text{Parameters} \approx 11.1\,\text{M}, \quad \text{FLOPs} \approx 28.7\,\text{G}
\end{equation}

This parameter count represents an optimal trade-off for the hardware constraints of the target system, offering a theoretical inference speed suitable for real-time monitoring ($>$20 FPS), while remaining significantly lighter than the YOLOv8m variant (25.9M parameters).

\subsection{Loss Function Topology}
The model minimizes a compound loss function $\mathcal{L}_{total}$ consisting of detection and classification components:
\begin{equation}
    \mathcal{L}_{total} = \lambda_{box}\mathcal{L}_{box} + \lambda_{cls}\mathcal{L}_{cls} + \lambda_{dfl}\mathcal{L}_{dfl}
\end{equation}
Where:
\begin{itemize}
    \item $\mathcal{L}_{box}$: The Complete Intersection over Union (CIoU) loss, which penalizes bounding box misalignment by simultaneously considering aspect ratio, centroid distance, and overlap area.
    \item $\mathcal{L}_{dfl}$: Distribution Focal Loss, used to refine the localization of box edges when defect boundaries are ambiguous, which is particularly relevant for thin stringing or gradually lifting warping edges.
\end{itemize}

\section{Dataset Engineering and Domain Adaptation}
\label{sec:dataset}

The synthetic generation pipeline described in Section~\ref{sec:synthetic_data} was refined iteratively throughout this project. This section focuses on the key improvements introduced in the final training iteration and on the overall composition of the hybrid dataset.

The most significant pipeline improvement was the replacement of the initial \textit{centroid alignment} strategy with \textbf{full-bed randomized placement}. In the original approach, defects were always composited at the geometric center of the printer bed background, scaled to occupy a large fraction of the bed area. While this produced clean, well-labelled samples, it introduced two biases: a positional bias, in which the model learned to expect defects in the center of the frame, and a scale bias, in which it learned to predict oversized bounding boxes that filled most of the frame. Neither reflected real-world behaviour, where failures can originate anywhere on the bed and where a localized defect occupies only a fraction of a large-format print. In the revised pipeline, both the position and the scale of each composited defect are randomized: placement spans the full bed area (with a 15\% edge margin), and scale is sampled uniformly between 10\% and 45\% of the bed width. This change substantially improved both the model's spatial generalization and the tightness of its predicted boxes.

A second, equally important change was the move from a purely synthetic dataset to a \textbf{hybrid corpus} that adds real defect imagery. For each class, approximately 150 synthetic composites are combined with approximately 150 real close-up defect images sourced directly from the Kaggle dataset, the latter labelled with a single near-full-frame bounding box ($0.5,\ 0.5,\ 0.85,\ 0.85$ in normalized coordinates) since the defect fills most of the close-up frame. This gives the model both the global spatial context of the synthetic views and the genuine texture and appearance of real defects, across the five trained classes (Spaghetti, Warping, Layer\_shifting, Stringing, Cracking) --- roughly 1{,}500 positive samples in total. Two classes, Warping and Cracking, are a deliberate exception: their only available real source images are the homogeneous gold-object laboratory shots discussed in Section~\ref{sec:results}, so they are kept synthetic-only to avoid reinforcing the colour bias. To these positives are added roughly 300 negative (defect-free) images --- 52 photographs of successful prints, around 250 clean webcam frames, and a small set of hard negatives --- giving a negative ratio of about 17\%. The complete corpus was divided into an 80/20 train/validation split, ensuring that the validation set contains samples the model never encountered during training.

\section{Training Protocol}
\label{sec:training}

The model was trained from the official \texttt{yolov8s.pt} pretrained base weights using Stochastic Gradient Descent (SGD) with a momentum of 0.937. Training ran for a maximum of 50 epochs with an early stopping patience of 15 epochs. To simulate the optical variability of a frame-mounted camera operating under real-world conditions, a comprehensive augmentation pipeline was applied:

\begin{table}[h]
\centering
\caption{Training Hyperparameters and Augmentation Settings}
\label{tab:hyperparams}
\begin{tabular}{|l|l|l|}
\hline
\textbf{Parameter} & \textbf{Value} & \textbf{Justification} \\ \hline
Input Resolution & $640 \times 640$ & YOLOv8 native input standard \\ \hline
Batch Size & 16 & VRAM optimization \\ \hline
Epochs & 50 (patience 15) & Early stopping to prevent overfitting \\ \hline
Mosaic & 0.5 & Simulates partial and zoom-out views \\ \hline
MixUp & 0.1 & Enhances robustness to overlapping defects \\ \hline
Rotation & $\pm 5^{\circ}$ & Simulates camera mounting variances \\ \hline
Horizontal Flip & 0.5 & Symmetry augmentation \\ \hline
HSV Hue & 0.015 & Simulates white balance variation \\ \hline
HSV Saturation & 0.4 & Simulates lighting intensity changes \\ \hline
HSV Value & 0.4 & Simulates exposure variation \\ \hline
\end{tabular}
\end{table}

The HSV augmentations deserve particular mention, as they were specifically motivated by the deployment context. A USB camera mounted inside a printer enclosure is exposed to variable illumination: the heated bed glows at high temperatures, ambient light changes throughout the day, and filament color varies between prints. Training with HSV perturbations encourages the model to rely on defect morphology rather than absolute brightness, improving robustness to these conditions without requiring additional data collection. It is worth noting in advance, however, that the modest default hue perturbation ($\text{hsv\_h}=0.015$) proved insufficient to fully break the colour--class correlation present in the Cracking and Warping source data; this limitation, and a more aggressive hue-randomisation strategy to address it, are analysed in Section~\ref{sec:results} and Chapter~\ref{chap:conclusion}.

\section{Experimental Results}
\label{sec:results}

\subsection{Quantitative Analysis}

The training process followed an iterative development approach common in applied machine learning research: an initial baseline run was first used to validate the end-to-end pipeline, confirming that the architecture, the synthetic data generation, and the augmentation strategy were producing meaningful learning signals. Subsequent iterations progressively refined the data composition and augmentation strategy. The model selected for deployment (\texttt{yolov8s\_improved\_v5}), trained and evaluated on a held-out 20\% validation split, achieved the following results:

\begin{table}[h]
\centering
\caption{Model Performance on Held-Out Validation Set (\texttt{yolov8s\_improved\_v5})}
\label{tab:results}
\begin{tabular}{|l|l|}
\hline
\textbf{Metric} & \textbf{Value} \\ \hline
mAP@50          & 0.970 \\ \hline
mAP@50--95      & 0.825 \\ \hline
Precision       & 0.969 \\ \hline
Recall          & 0.941 \\ \hline
\end{tabular}
\end{table}

These results are strong. A mAP@50 of 0.970 on a properly separated validation set indicates that the model has learned robust and generalizable features for the five trained defect classes across a variety of synthetic and real compositions. Classes with distinctive morphologies, such as \textit{Spaghetti} (chaotic filament tangles) and \textit{Layer Shifting} (staircase offsets), showed consistently high per-class precision. The bounding box regression loss decreased steadily over training, confirming that the model successfully learned to localize defects at the global field-of-view scale rather than relying solely on close-up textural cues. The relatively lower mAP@50--95 (0.825) reflects the looser localization that is expected for the deliberately large, near-full-frame bounding boxes used for the real defect samples, rather than a failure of classification.

\begin{figure}[h]
    \centering
    \includegraphics[width=0.8\textwidth]{images/results.png}
    \caption{Training metrics over the full training run. The steady decline in box loss and the plateau of mAP@50 indicate stable convergence without overfitting on the held-out validation set.}
    \label{fig:results}
\end{figure}

\subsection{Real-World Generalization and Open Challenges}

Despite strong validation metrics, real-world testing revealed an important limitation that is characteristic of models trained on narrowly-sourced data: the model's performance degrades on inputs that fall outside the training distribution. Testing on recorded video of actual print sessions showed encouraging results for visually salient defects such as Spaghetti and Layer Shifting, but two systematic failure modes emerged.

\textbf{Colour and domain bias.} The publicly available Cracking and Warping samples were all captured in the same laboratory setup --- small gold and brass test specimens on a white background, often accompanied by a yellow warning label. Because these classes lacked the diversity of the others, the model learned a spurious correlation between the \textit{colour and context} of the source objects and the defect label, rather than the defect morphology itself: golden or metallic prints tended to trigger false Cracking detections, and the characteristic lab framing tended to trigger false Warping detections. This is a textbook example of dataset bias, and it is the single most important open problem identified in this work.

\textbf{Temporal flicker.} On video, isolated frames occasionally produced spurious high-confidence detections, which would be undesirable in a system that can pause a multi-day print.

Two inference-side mitigations were applied. First, rather than a single global confidence threshold, \textit{per-class} thresholds were derived from the validation F1-curves and confusion matrix (Spaghetti $0.62$, Stringing $0.57$, Warping $0.55$, Cracking $0.52$, Layer\_shifting $0.50$), raising the bar most for the class the confusion matrix showed to be the most false-positive-prone (Spaghetti). Second, a temporal smoother requires a class to be detected in at least $5$ of the last $7$ frames before it is reported, suppressing single-frame flicker. These measures make the current model a usable and honest foundation for the monitoring system. They do not, however, remove the underlying bias; eliminating it requires either collecting colour-diverse real defect images or, as a zero-cost alternative explored as future work, randomising the hue of the extracted defect foregrounds during synthetic composition so that the model can no longer associate a specific colour with a specific defect. Addressing this distributional bias is identified as the highest-priority direction for future work (Chapter~\ref{chap:conclusion}).

% --------------------------------------------------
\chapter{Real-Time System Implementation}
\label{chap:implementation}

This chapter describes how the trained detector is embedded into a continuously running monitoring loop. The implementation follows the modular architecture introduced in Chapter~\ref{chap:architecture} and is designed from the outset for eventual deployment on a Raspberry Pi edge device. At the current stage it is fully implemented and validated in a \textit{software-in-the-loop} (mock) configuration, in which a recorded print video stands in for the live camera and a simulated printer interface stands in for the physical machine. The hardware-facing components are defined behind clean interfaces so that the transition to live operation requires no change to the control logic.

\section{System Integration}
The monitoring application is organised around three decoupled components, each hiding its implementation behind a narrow interface so that the mock and live variants are interchangeable:

\begin{itemize}
    \item \textbf{Camera} --- abstracts frame acquisition. In mock mode it reads frames from a recorded video file; in live mode it opens a USB webcam through its device index. Both paths expose the same frame-grabbing method returning a single BGR frame, so the rest of the system is agnostic to the source.
    \item \textbf{Detector} --- wraps the trained YOLOv8s model. It runs inference on an incoming frame, applies the per-class confidence filtering described in Chapter~\ref{chap:model_development}, and maintains the temporal persistence state used to debounce detections.
    \item \textbf{Printer Interface} --- abstracts communication with the printer firmware. In mock mode it returns safe status codes and logs the actions it would take; in live mode it is designed to issue HTTP requests to an OctoPrint or Moonraker server.
\end{itemize}

A central control loop --- the system entry point --- wires these components together: it acquires a frame, passes it to the detector, and, when a defect is confirmed, queries the printer state and issues a pause command. This separation of concerns means that validating the full pipeline against recorded footage exercises exactly the same control logic that will later run against live hardware; only the two leaf components, the camera source and the printer endpoint, are swapped.

\section{Inference Pipeline}
\textbf{Frame sampling.} Although the camera can deliver frames at its native rate, the loop deliberately samples at a low target rate of one frame per second. This choice is driven by the physics of the problem rather than by a hardware limitation: the targeted defects evolve over seconds to minutes, so sub-second sampling would yield highly redundant frames while needlessly consuming CPU on the edge device. Sampling at 1\,FPS leaves a generous per-frame time budget for inference and keeps the Raspberry Pi thermally and computationally comfortable.

\textbf{Per-frame inference.} Each sampled frame is passed to the YOLOv8s model. With approximately 11\,M parameters (Chapter~\ref{chap:model_development}), the network is light enough to run comfortably within the one-second budget, and an ONNX export of the trained weights is used to accelerate inference on the edge device. Because the sampling interval, and not the network, is the rate-limiting step, the real-time constraint is met by a wide margin.

\textbf{Temporal persistence.} A single-frame detection is never sufficient to act upon. The live detector applies a \textit{persistence filter}: a rolling counter is incremented on every frame in which a defect is present and reset to zero on any clean frame, and an alert is raised only once the counter reaches a configurable threshold (five consecutive frames by default). At 1\,FPS this means a defect must remain visible for roughly five continuous seconds before the system reacts. For the offline analysis of recorded video, a slightly stricter variant is used --- a sliding-window smoother that requires a class to appear in at least five of the last seven frames --- which tolerates brief dropouts while still suppressing isolated false positives. Both mechanisms trade a few seconds of detection latency, negligible against a multi-hour print, for a large reduction in spurious alerts.

\section{Alert and Response Mechanism}
When the persistence filter confirms a defect, the control loop initiates the response sequence. It first queries the printer interface to confirm that a job is actually in progress, and then issues a pause command; on success, the detector's persistence state is reset so that the system does not immediately re-trigger on the same, now-frozen, defect. Requiring confirmation across several consecutive frames before pausing is a deliberate safety design: pausing a multi-day print on the strength of a single false-positive frame would be unacceptable, so the system is biased toward acting only on persistent evidence.

The printer interface itself is built as a mode-switchable abstraction. In the \textbf{mock} mode used throughout development, the status-query and pause operations return safe values and log their actions, allowing the entire response path to be exercised without any risk to a physical machine. The \textbf{live} mode is designed around the REST APIs of the two most common open-source print controllers --- OctoPrint and Moonraker --- with the corresponding job-status and pause endpoints documented in the interface; connecting these calls to a physical printer is the immediate next step toward field deployment and is discussed as future work in Chapter~\ref{chap:conclusion}. A complementary user-notification channel --- for example an e-mail or push alert raised on detection --- is foreseen in the architecture as a dedicated notifier component and is likewise left for future implementation.

% --------------------------------------------------
\chapter{Experimental Results}
\label{chap:results}

\section{Testing Scenarios}
Normal printing vs anomalous cases.

\section{Results}
Performance of the system in detecting different defects.

\section{Discussion}
Limitations, false positives, and reliability of system.

% --------------------------------------------------
\chapter{Conclusion and Future Work}
\label{chap:conclusion}

\section{Conclusion}
This thesis set out to design and implement an AI-based visual monitoring system capable of detecting anomalies in real time during large-format FDM printing, with the concrete aim of safeguarding the multi-day prints used to manufacture humanoid-robot parts on an Elegoo OrangeStorm Giga. The work delivered four main contributions.

First, it addressed the scarcity of suitable training data through a \textbf{hybrid dataset methodology} that composites background-removed defect features onto real images of the target print bed, augmented with genuine defect close-ups and a body of defect-free negatives. This directly tackles the domain gap between the macro close-ups available publicly and the global field of view of a frame-mounted camera (Chapters~\ref{chap:model} and~\ref{chap:related}).

Second, it produced a \textbf{trained detector} --- the YOLOv8s model \texttt{yolov8s\_improved\_v5} --- that reaches a mAP@50 of 0.970 on a held-out validation set across five defect classes, confirming that the architecture and data strategy yield strong, generalisable features (Section~\ref{sec:results}).

Third, it implemented a \textbf{modular, edge-oriented real-time pipeline} --- acquisition, detection, temporal decision filtering, and printer actuation --- validated end-to-end in a software-in-the-loop configuration and designed for deployment on a Raspberry Pi mounted on the printer (Chapters~\ref{chap:architecture} and~\ref{chap:implementation}).

Fourth, and equally important, it provided an \textbf{honest characterisation of the system's principal limitation}: a colour and domain bias inherited from homogeneous source data, which causes false positives on certain colours and contexts (Section~\ref{sec:results}). Identifying this clearly, rather than masking it behind favourable validation numbers, is what makes the system a trustworthy foundation to build upon.

Taken together, these results demonstrate a viable, low-cost path toward autonomous monitoring of large-format printing. The remaining distance between strong validation performance and full real-world robustness defines the work that follows.

\section{Future Work}
The natural trajectory of this project is to evolve the current research prototype into a complete, deployable monitoring \textit{application} --- comparable in scope and polish to the integrated commercial print-monitoring solutions reviewed during this work, but open, self-hosted, and tailored to large-format production. The path toward that goal can be grouped into near-term engineering and longer-term research.

\subsection{Toward a Deployable Application}
\begin{itemize}
    \item \textbf{Close the control loop.} Complete the live OctoPrint/Moonraker integration so that the system issues real pause commands to the Elegoo OrangeStorm Giga, moving from the validated mock pipeline (Chapter~\ref{chap:implementation}) to true closed-loop operation.
    \item \textbf{Eliminate the colour and domain bias.} Collect colour-diverse real defect imagery for the affected classes, add metallic-but-clean prints as hard negatives, and randomise the hue of the extracted defect foregrounds during synthetic composition together with more aggressive HSV augmentation. This is the single highest-priority improvement, as it directly attacks the limitation identified in Section~\ref{sec:results}.
    \item \textbf{Optimise and benchmark on the edge device.} Profile the exported ONNX model on the Raspberry Pi, apply quantisation if required, and measure sustained frame rate and thermal behaviour to confirm real-time operation under deployment conditions.
    \item \textbf{Productise the prototype.} Package the system as a persistent background service with a user-facing web dashboard --- live camera view, real-time alerts, and a history of detected events --- together with configuration controls and notification channels such as e-mail or push messages. This is the step that turns a collection of scripts into the integrated application that is the ultimate goal of the project.
    \item \textbf{Continuous, self-improving operation.} Activate passive data collection during real prints so that the system gathers genuine, environment-specific samples as it runs, enabling periodic retraining in an active-learning loop that improves accuracy the longer the system is deployed.
    \item \textbf{Broader defect coverage.} Extend the model to the two defect classes already part of the design taxonomy but not yet trained --- Off-Platform and Blob of Death (Chapter~\ref{chap:defects}) --- and, in time, to subtler faults such as under- and over-extrusion.
\end{itemize}

\subsection{Longer-Term Extensions}
\begin{itemize}
    \item \textbf{Multi-camera coverage.} A single camera cannot fully resolve a bed approaching a metre across; multiple synchronised views would eliminate blind spots and improve the detection of defects near the edges.
    \item \textbf{Thermal imaging and multi-sensor fusion.} Combining the visual stream with a thermal camera and with the process signals discussed in Chapter~\ref{chap:related} would allow the system to detect failure modes that are not visually obvious, such as incipient warping driven by uneven cooling.
    \item \textbf{Automatic parameter optimisation.} Beyond pausing the print, a mature system could close the loop more tightly by adjusting print parameters in response to detected anomalies, in the spirit of the autonomous optimisation demonstrated by Zhang et al.\ \cite{zhang2025} --- turning the monitor from a safety net into an active quality-control agent.
\end{itemize}


\backmatter



\begin{thebibliography}{99}



\bibitem{armin2025}

E. Armin et al., ``Defect detection in 3D printing: a review of image processing and machine vision techniques,'' \textit{The International Journal of Advanced Manufacturing Technology}, vol.~140, pp.~2103--2128, 2025.



\bibitem{zhang2025}

H. Zhang et al., ``Autonomous printing process optimisation and in-situ anomaly detection in fused deposition modelling using an integrated data-driven approach,'' \textit{Virtual and Physical Prototyping}, vol.~20, no.~1, e2545523, 2025.



\bibitem{cenedese2022}

E. Cenedese, ``3D Printing Anomaly Detection: Implementation of a ML system using YOLOv5 and EfficientNet-Lite,'' M.S. thesis, Politecnico di Torino, 2022.



\bibitem{cascio2017}

M. Cascio, ``Un Approccio Innovativo basato su Deep Learning per la Classificazione di Human Action,'' M.S. thesis, Sapienza Università di Roma, 2017.



\bibitem{fagioli2019}

A. Fagioli, ``Non-Acted Body Affect Recognition,'' M.S. thesis, Sapienza Università di Roma, 2019.



\end{thebibliography}



\end{document}
